\section{Limitations and caveats}
\label{sec:limits}

\begin{itemize}[leftmargin=*,itemsep=2pt]
\item \textbf{Timing noise and sharing.} Timings come from fresh processes on shared cluster nodes. GPU nodes were not exclusive during the main matrix, so the other GPU of a two-GPU node (and the filesystem) may have been in use; medians over five to nine samples bound this (the spread of the individual cases is visible in Figures~\ref{fig:ratio} and~\ref{fig:violin}), but ratios within a few percent of one should be read as ``no difference''. Disk-tier timings depend on the page cache and the shared filesystem: two alternating runs of identical code on one node differed by up to 12\%, and repeated runs of Taxol on the V100 ranged from 239 to 280\,s. No general I/O performance claim is made; the disk-tier ratios refer to the specific runs and file system.
\item \textbf{Small systems.} Many core cases finish in one to a few seconds, where start-up costs, kernel launch latency and the CPU-side set-up dominate; ratios near one for these are expected and do not measure the numerical kernels.
\item \textbf{Memory figures} are the peak of requested live CuPy storage in a separate instrumented process. They exclude the CuPy pool overhead, cuTENSOR and cuBLAS workspaces and native allocations, so the real device footprint is somewhat higher, and the P100-16 rows that report \emph{resource limited} without a measurement are forecasts made by the scheduler before launch (they are labelled in the data).
\item \textbf{Capacity runs} have no independent reference; they are checked only for finite values and SCF convergence. The largest system with an accuracy-qualified comparison to CPU PySCF is Taxol.
\item \textbf{Registry amendments} (Section~\ref{sec:deviations}) change the numerical protocol relative to the original plan. In particular the GPU TDA calculations run at a tighter solver tolerance than the TDDFT calculations, so the TDA and TDDFT times must not be compared with each other; and the CPU Hessian timing is a finite-difference route, not PySCF's analytic Hessian.
\item \textbf{Non-converged and ill-conditioned cases.} Four open-shell cases do not converge on CPU and are excluded; two of the three 48-thread CPU repetitions of NO/B3LYP differ from the 24-thread reference by $1.1\times10^{-8}$ and $1.3\times10^{-8}$~Eh (gate $10^{-8}$) and are kept as accuracy failures; two of five benzene TDA and TDDFT roots are flagged non-converged by the GPU solver at the declared tolerance for both code versions, although the excitation energies agree with CPU to $1.2\times10^{-9}$~Eh.
\item \textbf{The original code fails the CPU gradient gate on Taxol} for HF, B3LYP and wB97X (maximum component error about $2\times10^{-5}$~Eh/Bohr), on all cards. Those cases have no equal-accuracy ratio.
\item \textbf{Runner labelling.} Plain CUDA out-of-memory errors (\texttt{cudaErrorMemoryAllocation}) are recorded by the runner as \emph{infrastructure failures}; they are genuine resource limits and are discussed as such. Pascal GPUs also fail CCSD with a native kernel launch error whose cause was not diagnosed.
\item \textbf{Audit tolerance.} The suite's audit demands exact equality of recomputed comparison statistics; last-digit differences between processes (multi-threaded reductions) required a wrapper that compares floating-point values to a relative $10^{-9}$ and everything else exactly. Pre-launch admission rejections are accepted without a request file only when they carry their forecast.
\item \textbf{Code-review findings left unchanged:} the 128\,MiB default staging buffer (the campaign sets it explicitly), the GGA density shortcut (empirically within every gate, but its sum-of-squares form is no longer exactly non-negative), the \texttt{strict\_grid\_order} compatibility flag, the missing early guard for disk tensors in the old DF-MP2 module, and a duplicated staging-cap check.
\item \textbf{CPU baseline.} The CPU is not optimized in this study. A 48-thread whole-node lane (the whole-node default) and a 24-thread, one-BLAS-thread lane are reported; the fair baseline for each case is the faster of the lanes. Other CPU settings (other thread counts, different BLAS) were not scanned for every case.
\end{itemize}
